% =============================================================
%  Chapter — Control & Protection (Module 10)
%  Figures in ./fig  (EMF converted to cropped PDF, PNG copied).
% =============================================================
\newcommand{\cpfig}[1]{chapters/c-and-p/fig/#1}

\sectionslide{Control \& Protection}{Converter control and dynamic behaviour}

\begin{frame}{What can a converter control?}
  \begin{columns}[c,onlytextwidth]
    \begin{column}{0.54\textwidth}
      \textbf{Control values of the fast control loop:}
      \begin{itemize}
        \setlength{\itemsep}{0.25em}
        \item DC voltage
        \item DC current
        \item AC current
          \begin{itemize}
            \item Positive-phase active current $i_{p,psf}$
            \item Positive-phase reactive current $i_{q,psf}$
            \item Negative-phase active current $i_{p,nsf}$
            \item Negative-phase reactive current $i_{q,nsf}$
          \end{itemize}
        \item AC voltage
        \item AC frequency
        \item Converter energy
      \end{itemize}
    \end{column}
    \begin{column}{0.42\textwidth}
      \centering
      \topoimg[width=\linewidth,height=0.66\textheight,keepaspectratio]{\cpfig{image4}}
    \end{column}
  \end{columns}
\end{frame}

\begin{frame}{Energy flow}
  \begin{columns}[c,onlytextwidth]
    \begin{column}{0.54\textwidth}
      \begin{itemize}
        \setlength{\itemsep}{0.5em}
        \item The converter has only limited energy storage, so in steady
              state the converter energy must be balanced over one period.
        \item Energy balancing can be done through the AC side or through
              the DC side.
        \item One side defines the power flow, the other is responsible for
              the balancing.
      \end{itemize}
    \end{column}
    \begin{column}{0.42\textwidth}
      \centering
      \topoimg[width=\linewidth,height=0.32\textheight,keepaspectratio]{\cpfig{image5}}\\[0.6em]
      \topoimg[width=\linewidth,height=0.32\textheight,keepaspectratio]{\cpfig{image6}}
    \end{column}
  \end{columns}
\end{frame}

\begin{frame}{Basic control principle of the HVDC converter}
  \begin{columns}[c,onlytextwidth]
    \begin{column}{0.50\textwidth}
      \textbf{Three major control loops:}
      \begin{itemize}
        \setlength{\itemsep}{0.25em}
        \item AC loop
        \item DC loop
        \item Inner loop
      \end{itemize}
      \vspace{0.5em}
      \textbf{Impedances of the control loop:}
      \begin{itemize}
        \setlength{\itemsep}{0.2em}
        \item AC converter impedance
        \item DC converter impedance
        \item AC grid impedance
        \item DC line impedance
        \item Converter internal impedance
      \end{itemize}
    \end{column}
    \begin{column}{0.46\textwidth}
      \centering
      \topoimg[width=\linewidth,height=0.68\textheight,keepaspectratio]{\cpfig{image7}}
    \end{column}
  \end{columns}
\end{frame}

\begin{frame}{AC current control}
  \begin{columns}[c,onlytextwidth]
    \begin{column}{0.50\textwidth}
      \textbf{Control input / reference:} $\omega_{conv}$, $i_p$\,/\,$P$\\[0.4em]
      \textbf{Control output:} $U_{dc}$
      \vspace{0.6em}

      \textbf{Scenario}
      \begin{itemize}
        \setlength{\itemsep}{0.3em}
        \item AC power is defined by the operator.
        \item DC voltage controls the power flow on the DC side to balance
              the converter energy.
      \end{itemize}
    \end{column}
    \begin{column}{0.46\textwidth}
      \centering
      \topoimg[width=\linewidth,height=0.66\textheight,keepaspectratio]{\cpfig{image8}}
    \end{column}
  \end{columns}
\end{frame}

\begin{frame}{AC voltage control}
  \begin{columns}[c,onlytextwidth]
    \begin{column}{0.50\textwidth}
      \textbf{Control input / reference:} $\omega_{conv}$, $U_{AC}$, $f_{AC}$\\[0.4em]
      \textbf{Control output:} $U_{dc}$
      \vspace{0.6em}

      \textbf{Scenario}
      \begin{itemize}
        \setlength{\itemsep}{0.3em}
        \item AC power is defined by the AC load.
        \item DC voltage controls the power flow on the DC side to balance
              the converter energy.
      \end{itemize}
    \end{column}
    \begin{column}{0.46\textwidth}
      \centering
      \topoimg[width=\linewidth,height=0.66\textheight,keepaspectratio]{\cpfig{image9}}
    \end{column}
  \end{columns}
\end{frame}

\begin{frame}{DC voltage control}
  \begin{columns}[c,onlytextwidth]
    \begin{column}{0.50\textwidth}
      \textbf{Control input / reference:} $\omega_{conv}$, $U_{DC}$\\[0.4em]
      \textbf{Control output:} $i_p$
      \vspace{0.6em}

      \textbf{Scenario}
      \begin{itemize}
        \setlength{\itemsep}{0.3em}
        \item The station defines the voltage on the DC line.
        \item AC current (and therefore AC power) balances the converter
              energy.
      \end{itemize}
    \end{column}
    \begin{column}{0.46\textwidth}
      \centering
      \topoimg[width=\linewidth,height=0.66\textheight,keepaspectratio]{\cpfig{image10}}
    \end{column}
  \end{columns}
\end{frame}


\begin{frame}{Reactive power control}
  \begin{columns}[c,onlytextwidth]
    \begin{column}{0.52\textwidth}
      \begin{itemize}
        \setlength{\itemsep}{0.5em}
        \item In control modes without AC voltage and frequency control, the
              reactive power can be controlled independently of the energy
              balancing.
        \item Reactive power does not require DC power flow and can operate
              without the DC line connected.
      \end{itemize}
    \end{column}
    \begin{column}{0.44\textwidth}
      \centering
      \topoimg[width=\linewidth,height=0.66\textheight,keepaspectratio]{\cpfig{image8}}
    \end{column}
  \end{columns}
\end{frame}


\begin{frame}{AC control - Current source vs. Voltage source}
  \begin{columns}[c,onlytextwidth]
    \begin{column}{0.52\textwidth}
      \centering
      \textbf{AC current source}\par
      \vspace{1em}
      \topoimg[width=\linewidth,height=0.5\textheight,keepaspectratio]{\cpfig{image12}}
    \end{column}
    \begin{column}{0.44\textwidth}
      \centering
      \textbf{AC voltage source}\par
      \vspace{1em}
      \topoimg[width=\linewidth,height=0.5\textheight,keepaspectratio]{\cpfig{image11}}
    \end{column}
  \end{columns}
\end{frame}


\begin{frame}{AC control - Current source vs. Voltage source}
  \centering
  \renewcommand{\arraystretch}{1.35}
  \begin{tabular}{@{}p{5.6cm} p{4cm} p{3.5cm}@{}}
     & \textbf{AC current source} & \textbf{AC voltage source} \\
    \hline
     & Grid Following (GFL)\newline P/Q control mode & V/f or \newline Grid Forming (GFM) \\
    \textbf{Primary control parameters}   & Active current\newline Reactive current & AC frequency\newline AC voltage \\
    \textbf{Secondary control parameters} & Active power\newline Reactive power\newline AC voltage & \\
    \textbf{Grid application}             & SCL $\ge 2$ & SCL $<2$ \\
    \hline
  \end{tabular}
\end{frame}


\begin{frame}{Grid-forming control concept}
  \centering
  \topoimg[width=\linewidth,height=0.72\textheight,keepaspectratio]{\cpfig{image13}}\\[0.3em]
  {\color{atGray}\footnotesize Grid-forming converters - advantages of grid-forming control in HVDC and FACTS applications\\
  (white paper, Siemens Energy).}
\end{frame}

\begin{frame}{Steady-state power transmission}
  \begin{columns}[c,onlytextwidth]
    \begin{column}{0.52\textwidth}
      \begin{itemize}
        \setlength{\itemsep}{0.4em}
        \item One station controls DC voltage.
        \item One station controls active power.
        \item Two active-power control modes are possible:
          \begin{itemize}
            \item P/Q control mode
            \item U/f control mode
          \end{itemize}
      \end{itemize}
    \end{column}
    \begin{column}{0.44\textwidth}
      \centering
      \topoimg[width=\linewidth,height=0.32\textheight,keepaspectratio]{\cpfig{image14}}\\[0.6em]
      \topoimg[width=\linewidth,height=0.32\textheight,keepaspectratio]{\cpfig{image15}}
    \end{column}
  \end{columns}
\end{frame}

\begin{frame}{Steady-state control principle}
  \begin{columns}[T,onlytextwidth]
    \begin{column}{0.48\textwidth}
      \textbf{Standard operation}
      \begin{itemize}
        \setlength{\itemsep}{0.25em}
        \item The station that sees the highest DC voltage controls the DC voltage.
        \item Active power can be defined at any station.
      \end{itemize}
    \end{column}
    \begin{column}{0.48\textwidth}
      \topoimg[width=0.9\linewidth,height=0.28\textheight,keepaspectratio]{\cpfig{image17}}
    \end{column}
  \end{columns}
\end{frame}

\begin{frame}{Steady-state control principle}
  \begin{columns}[T,onlytextwidth]
    \begin{column}{0.48\textwidth}
      \textbf{Wind-farm connection}
      \begin{itemize}
        \setlength{\itemsep}{0.25em}
        \item Rectifier receives power coming from the WTG.
        \item Inverter must consider the DC line voltage drop.
      \end{itemize}
    \end{column}
    \begin{column}{0.48\textwidth}
      \topoimg[width=0.9\linewidth,height=0.28\textheight,keepaspectratio]{\cpfig{image18}}
    \end{column}
  \end{columns}
\end{frame}


\begin{frame}{Steady-state control principle}
  \begin{columns}[T,onlytextwidth]
    \begin{column}{0.48\textwidth}
      \textbf{Power supply to islanded network}
      \begin{itemize}
        \setlength{\itemsep}{0.25em}
        \item Inverter provides power required by the loads
        \item Rectifier controls DC voltage by providing the required active power to the DC side
      \end{itemize}
    \end{column}
    \begin{column}{0.48\textwidth}
      \topoimg[width=0.9\linewidth,height=0.28\textheight,keepaspectratio]{\cpfig{image19}}
    \end{column}
  \end{columns}
\end{frame}


\begin{frame}{Dynamic analysis}
  \begin{columns}[c,onlytextwidth]
    \begin{column}{0.4\textwidth}
      \begin{itemize}
        \setlength{\itemsep}{0.4em}
        \item Dynamic analysis covers converter behaviour during AC grid
              faults (on the network side of the converter transformer).
        \item AC grid faults are very common, so the converter must stay
              connected and provide supporting functions.
      \end{itemize}
      \vspace{0.3em}
    \end{column}
    \begin{column}{0.6\textwidth}
      \centering
      \topoimg[width=\linewidth,height=0.66\textheight,keepaspectratio]{\cpfig{image20}}
      {\color{atGray}\footnotesize CIGR\'E TB 832}
    \end{column}
  \end{columns}
\end{frame}

\begin{frame}{Dynamic analysis: 3ph-to-GND at the inverter}
  \centering
  \begin{columns}[c,onlytextwidth]
    \begin{column}{0.32\textwidth}\centering
      \topoimg[width=\linewidth,height=0.8\textheight,keepaspectratio]{\cpfig{image21}}\end{column}
    \begin{column}{0.32\textwidth}\centering
      \topoimg[width=\linewidth,height=0.8\textheight,keepaspectratio]{\cpfig{image22}}\end{column}
    \begin{column}{0.32\textwidth}\centering
      \topoimg[width=\linewidth,height=0.8\textheight,keepaspectratio]{\cpfig{image23}}\end{column}
  \end{columns}
  \vspace{0.3em}
  {\color{atGray}\footnotesize CIGR\'E TB 832}
\end{frame}

\begin{frame}{Dynamic analysis: 3ph-to-GND at the rectifier}
  \centering
  \begin{columns}[c,onlytextwidth]
    \begin{column}{0.32\textwidth}\centering
      \topoimg[width=\linewidth,height=0.8\textheight,keepaspectratio]{\cpfig{image24}}\end{column}
    \begin{column}{0.32\textwidth}\centering
      \topoimg[width=\linewidth,height=0.8\textheight,keepaspectratio]{\cpfig{image25}}\end{column}
    \begin{column}{0.32\textwidth}\centering
      \topoimg[width=\linewidth,height=0.8\textheight,keepaspectratio]{\cpfig{image26}}\end{column}
  \end{columns}
  \vspace{0.3em}
  {\color{atGray}\footnotesize CIGR\'E TB 832.}
\end{frame}

\begin{frame}{Control hierarchy}
  \begin{columns}[c,onlytextwidth]
    \begin{column}{0.48\textwidth}
      \centering
      \topoimg[width=\linewidth,height=0.66\textheight,keepaspectratio]{\cpfig{image27}}
    \end{column}
    \begin{column}{0.48\textwidth}
      \centering
      \topoimg[width=\linewidth,height=0.66\textheight,keepaspectratio]{\cpfig{image28}}
    \end{column}
  \end{columns}
  \vspace{0.3em}
  \centering
  {\color{atGray}\footnotesize CLC/TS 50654-1:2020 - HVDC control}
\end{frame}

\begin{frame}{Control requirements}
  \begin{columns}[c,onlytextwidth]
    \begin{column}{0.42\textwidth}
      \textbf{Input parameters for the control:}
      \begin{itemize}
        \setlength{\itemsep}{0.2em}
        \item Voltage measurement
        \item Current measurement
      \end{itemize}
      \vspace{0.4em}
      \textbf{Control functions:}
      \begin{itemize}
        \setlength{\itemsep}{0.2em}
        \item Current control
        \item Voltage control
        \item Reference voltage generation for six voltage sources
        \item Voltage synthesis by generating switching pulses for individual
              modules
      \end{itemize}
    \end{column}
    \begin{column}{0.58\textwidth}
      \centering
      \topoimg[width=\linewidth,height=0.66\textheight,keepaspectratio]{\cpfig{image29}}
    \end{column}
  \end{columns}
\end{frame}
